\chapter{Introduction}

\projectquote{``Everything should be made as simple as possible, but no simpler.''}

This template is designed for longer mathematics or physics documents that mix
structured exposition, reusable theorem environments, figures, exercises and
bibliography support.

\section{How to use it}

Start by editing \texttt{metadata.tex}, then replace these starter chapters with
your own material. The package setup is centralized in
\texttt{style/project.sty}, while reusable macros live in
\texttt{style/commands.tex}.

\defn{Template idea}{
    A good project template keeps recurring setup in one place so that future
    documents can focus on content rather than package cleanup.
}

Every \verb|\defn| title is added to the ``Index of Definitions'' at the end of
the document automatically. Use \verb|\defn[key]{Title}{...}| to index a
different key, or \verb|\defn[-]{Title}{...}| to skip the entry.

\rmkb{
If a future project needs fewer features, switch them off through the package
options in \texttt{main.tex} rather than letting the preamble grow ad hoc.
}

\citep{HaiounDefiningextendedTQFTshandleattachments2026}
\citep{AmbrosinoMonodromydefectsChernSimonstheoryHolography2026}
